\documentclass[../../../include/open-logic-section]{subfiles}

\begin{document}

\olfileid{sth}{story}{blv}

\olsection{પરિશિષ્ટ: ફ્રેગેનો મૂળભૂત નિયમ V}

\olref[rus]{sec}માં આપણે સમજાવ્યું હતું કે રસેલે પોતાનો
વિરોધાભાસ ફ્રેગેના \emph{Grundgesetze}માં રૂપરેખા આપેલી પદ્ધતિ
માટેની સમસ્યા તરીકે રજૂ કર્યો હતો. ફ્રેગેની પદ્ધતિમાં સહજ ગણરચનાનું
સીધું પ્રરૂપ નહોતું. તેથી આ પરિશિષ્ટમાં તેમની પદ્ધતિમાં
\emph{શું હતું}, તેનો સહજ ગણરચના સાથે શું સંબંધ છે અને રસેલના
વિરોધાભાસ સાથે શું સંબંધ છે તે બહુ ટૂંકમાં સમજાવીશું.

ફ્રેગેની પદ્ધતિ \emph{દ્વિતીય-ક્રમની} છે અને તે
\emph{સંકલ્પનાના વ્યાપ}ના ખ્યાલને વ્યક્ત કરવા માટે રચાઈ હતી.
\footnote{ચોક્કસ રીતે કહીએ તો ફ્રેગે વધુ સામાન્ય ખ્યાલને
ઔપચારિક બનાવવાનો પ્રયાસ કરે છે: વિધેયનો ``મૂલ્યવ્યાપ''.
સંકલ્પનાઓના વ્યાપો આ વધુ સામાન્ય ખ્યાલના વિશેષ કિસ્સા છે.
વિગતો માટે \citet[pp.\ 8--9]{Heck2012} જુઓ.}
ફ્રેગેથી પ્રેરિત સંજ્ઞા વાપરીને
\emph{સંકલ્પના~$F$ના વ્યાપ} માટે $\fregeext{x}{F(x)}$ લખીશું.
આ એવી રીત છે જે \emph{વિધેયધર્મ} ``$F$''ને લે છે અને તેને
(પ્રથમ-ક્રમના) \emph{પદ} ``$\fregeext{x}{F(x)}$''માં ફેરવે છે.
આ રીત વાપરીને ફ્રેગેએ સભ્યતાની નીચેની \emph{વ્યાખ્યા} આપી:
\[
a \in b =_\text{df} \exists G(b = \fregeext{x}{G(x)} \land Ga)
\]
સ્થૂળ રીતે: $a \in b$ ત્યારે અને માત્ર ત્યારે જ્યારે $a$ એવી
સંકલ્પના હેઠળ આવે જેનો વ્યાપ $b$ છે. (નોંધો કે પરિમાણક
``$\exists G$'' દ્વિતીય-ક્રમનો છે.) ફ્રેગેએ
\emph{મૂળભૂત નિયમ V} તરીકે ઓળખાતો નીચેનો સિદ્ધાંત પણ સ્વીકાર્યો:
$$\fregeext{x}{F(x)} = \fregeext{x}{G(x)} \liff \forall x (Fx \liff Gx)$$
સ્થૂળ રીતે: સંકલ્પનાઓના વ્યાપો અભિન્ન હોય છે ત્યારે અને માત્ર
ત્યારે જ્યારે તે સંકલ્પનાઓ સમાન વ્યાપવાળી હોય. (ફરીથી, ``$F$''
અને ``$G$'' બંને વિધેયધર્મના સ્થાને છે.) હવે એક સરળ સિદ્ધાંત
સભ્યતાને ગુણધર્મ સંતોષવા સાથે જોડે છે:

\begin{lem}[\emph{Grundgesetze}માં]\ollabel{lem:Fregeextensions}
$\forall F \forall a(a \in \fregeext{x}{F(x)} \liff Fa)$
\end{lem}

\begin{proof}
$F$ અને $a$ નક્કી કરો. હવે $a \in \fregeext{x}{F(x)}$ ત્યારે અને
માત્ર ત્યારે જ્યારે $\exists G(\fregeext{x}{F(x)}
= \fregeext{x}{G(x)} \land Ga)$ (સભ્યતાની વ્યાખ્યાથી), ત્યારે અને
માત્ર ત્યારે જ્યારે $\exists G(\forall x(Fx \liff Gx) \land Ga)$
(મૂળભૂત નિયમ Vથી), ત્યારે અને માત્ર ત્યારે જ્યારે $Fa$
(પ્રાથમિક દ્વિતીય-ક્રમ તર્કશાસ્ત્રથી).
\end{proof}

અને આમાંથી લગભગ તરત જ સહજ ગણરચના મળે છે:

\begin{lem}[\emph{Grundgesetze}માં.]
$\forall F \exists s \forall a (a \in s \liff Fa)$
\end{lem}

\begin{proof}
$F$ નક્કી કરો; હવે \olref{lem:Fregeextensions} પરથી
$\forall a (a \in
\fregeext{x}{F(x)} \liff Fa)$ મળે છે; તેથી અસ્તિત્વવાચક
સામાન્યીકરણ વડે $\exists s\forall a(a \in s \liff Fa)$ મળે છે.
$F$ મનસ્વી હોવાથી પરિણામ અનુસરે છે.
\end{proof}

$F$ને $\forall x(Fx \liff x \notin x)$ પ્રમાણે લઈએ તો રસેલનો
વિરોધાભાસ મળે છે.

\end{document}
